\chapter{Upgrade and migration}
\label{ch:upgrade}

Routine upgrades are designed to be safe. The exception is the
1.0.x to 1.1.0 jump, which is a one-time breaking change documented
below.

\section{Routine upgrades (1.y to 1.z)}

\begin{enumerate}
  \item Read the release notes for every version between your
        current and the target version.
  \item Take a fresh \textbf{database backup} (chapter
        \ref{ch:backup}).
  \item Update the image tag in \texttt{docker-compose.yml} (or
        repull \texttt{:latest} if you track that tag).
  \item Run \texttt{docker compose up -d}. The new container boots,
        runs any new migrations against the existing database, and
        starts serving requests.
  \item Watch the logs for migration errors:
        \texttt{docker compose logs mybibli | grep -i migrate}.
\end{enumerate}

Migrations are append-only — schema changes are forward
compatible. A 1.4 install upgraded to 1.5 keeps every row from 1.4
and applies new tables / columns in place.

\begin{tip}
For homelab installs, the \texttt{:latest} Docker tag is a
reasonable default. For production-style setups (a community
library, an association), pin to an explicit version tag so an
unattended pull does not bring in a release you have not read about.
\end{tip}

\section{The 1.0.x to 1.1.0 jump}\label{sec:upgrade-1.1.0}

\begin{warning}
\textbf{This is a one-time breaking change.}\index{1.1.0 upgrade}
1.0.x installs are subject to the issue documented in
\href{https://github.com/guycorbaz/mybibli/issues/173}{\#173}:
the seed migrations create the credentials \texttt{admin/admin}
and \texttt{librarian/librarian} on every fresh install, including
production, and silently bypass the first-launch setup wizard.
\end{warning}

The recommended path depends on whether your 1.0.x install holds
real data:

\subsection*{If your 1.0.x install is empty (no titles cataloged)}

The simplest path: wipe the database and reinstall on 1.1.0+.

\begin{enumerate}
  \item Stop the stack: \texttt{docker compose down}.
  \item Remove the database volume:
        \texttt{docker volume rm mybibli\_mybibli\_db\_data}
        (or the equivalent volume name on your host —
        \texttt{docker volume ls} lists them).
  \item Pull the 1.1.0 image:
        \texttt{docker compose pull}.
  \item Update the image tag in \texttt{docker-compose.yml} if you
        pinned a previous version.
  \item Start: \texttt{docker compose up -d}.
  \item Open \texttt{/setup} and walk through the wizard.
\end{enumerate}

\subsection*{If your 1.0.x install holds real data}

\begin{enumerate}
  \item Take a database backup (chapter \ref{ch:backup}).
  \item Take a covers + \texttt{.env} backup.
  \item Log in as the seeded admin and rotate the password
        \emph{before} you upgrade. The 1.1.0 upgrade flow itself is
        safe, but rotation closes the only window where the default
        credentials could leak.
  \item Update the image tag in \texttt{docker-compose.yml} and run
        \texttt{docker compose up -d}.
  \item The 1.1.0 boot detects the existing admin and does
        \emph{not} re-enable the wizard — the rotated admin remains
        the active admin.
\end{enumerate}

\subsection*{The new \texttt{MYBIBLI\_SEED\_DEV\_USERS} variable}

Starting with 1.1.0, the seed migration is gated by
\texttt{MYBIBLI\_SEED\_DEV\_USERS}\index{MYBIBLI\_SEED\_DEV\_USERS}.
The default is \texttt{0} (no seed). Only set it to \texttt{1} for
local-development workflows or automated tests — never in
production. The first-launch wizard becomes the canonical
onboarding path.

\section{Reading the release notes}

Every mybibli release publishes a GitHub Release page with:

\begin{itemize}
  \item a one-paragraph summary of the change set;
  \item the migration impact (none / additive / breaking);
  \item a link to the diff against the previous tag;
  \item attached PDFs of this manual (English + French) once the
        release-time PDF build is wired into CI.
\end{itemize}

Read the notes before upgrading. Breaking changes are tagged with a
prominent \emph{Breaking} label.

\section{Rollback}

Rollback is \emph{not} a first-class flow. Schema migrations move
forward only — there is no \texttt{down.sql}. The supported recovery
strategy is:

\begin{enumerate}
  \item Restore the pre-upgrade database backup
        (chapter~\ref{ch:backup}).
  \item Pin the image tag back to the prior version.
  \item Start the stack.
\end{enumerate}

This works because the backup is consistent with the prior schema.
Trying to point an older image at a newly-migrated database will
fail to boot.
